\scanpage{15}
\deleted{We say}\inserted{Let} $T$ be a neighborhood
\deleted{tubular}\inserted{conical} of $Y\subset X$
in $X$, put
\[
(11)\qquad X^*=X\setminus\mathring T=(X\setminus T)\cup\dot T
\]
it is a closed subset of $X$. \deleted{We say that}

\uline{Proposition\unsure{.}} Equivalent conditions (where $\lambda\in]0,1[$ given)

a) $\dot T$ is a boundary-subset of $X^*=X\setminus\mathring T$

b) There exists a homeomorphism of $X$
with itself, transforming $T,\dot T$ into
% crops/150_scan015_P04.png: the crossed lambda-like group may absorb a parenthesis; no separate live opening parenthesis is visible.
$\deleted{\unsure{\lambda}}T_\lambda,\dot T_\lambda)$

b') As in b, but assuming moreover
that the homeomorphism \unsure{induces|induces} the identity
on $T_{\lambda'}$ (where \deleted{$0<\lambda'<\lambda$}
$\lambda'$ is given, $0<\lambda'<\lambda$).

c) There exists a \deleted{tubular} conical neighborhood $T'$
of $Y$ in $X$, \deleted{such that and \unsure{a|\alpha}} such that $T=\lambda T'$,
the isomorphism being compatible with the
conical neighborhood structures.

We then say that $T$ is a \deleted{tubular} conical neighborhood
\uline{strict} in $X$

\margin{It is better to define conical neighborhoods
\unsure{strict|strict} as \unsure{some} \uline{germs} of a neighborhood
larger $(1+\varepsilon)T\unsure{-R_{>0}|-\text{then}}$
$(\varepsilon>0)$}

\uline{Corollary} When $T$ is a strict \deleted{tubular} conical neighborhood in $X$,
$\dot T$ is the boundary of $T$ in $X$.

\uline{Corollary} If \deleted{\unsure{it is|one}}
% crops/150_scan015_P11.png: additional triangular ink has an A-like outline; token/state and possible insertion-caret function remain uncertain.
\unsure{\deleted{$A$}|$A$}\inserted{$T$} neighborhood
\deleted{tubular}\inserted{conical} of $Y$
in $X$, the $T_\lambda$ $(\lambda\in]0,1[)$ are neighborhoods
\deleted{tubular}\inserted{conical} strict
\margin{14}
